\subsection{准希尔伯特空间的张量积}

设 \(E_{1}\) 与 \(E_{2}\) 为两个准希尔伯特空间，\(F = E_{1} \otimes E_{2}\) 为向量空间 \(E_{1}\) 与 \(E_{2}\) 的张量积。设 \(x_{1} \in E_{1}\) 且 \(x_{2} \in E_{2}\)；由于映射 \((y_{1}, y_{2}) \mapsto \langle x_{1}|y_{1} \rangle \langle x_{2}|y_{2} \rangle\)（从 \(E_{1} \times E_{2}\) 到 K 中）是双线性的，故存在线性型 \(\phi_{x_{1},x_{2}}\)（在 \(E_{1} \otimes E_{2}\) 上），使得

\[
\phi_ {x _ {1}, x _ {2}} (y _ {1} \otimes y _ {2}) = \langle x _ {1} | y _ {1} \rangle \langle x _ {2} | y _ {2} \rangle\tag{1}
\]

其中 \(y_{1} \in \mathrm{E}_{1}\)，\(y_{2} \in \mathrm{E}_{2}\)。设 \(z \in \mathrm{F}\)。映射 \((x_{1}, x_{2}) \mapsto \overline{\phi_{x_{1},x_{2}}(z)}\)（从 \(\mathrm{E}_1 \times \mathrm{E}_2\) 到 \(\mathbf{K}\) 中）是双线性的；把 \(z\) 写成形式 \(z = \sum_{i=1}^{n} y_{i,1} \otimes y_{i,2}\)（其中 \(y_{i,1} \in \mathrm{E}_1\)、\(y_{i,2} \in \mathrm{E}_2\)，\(1 \leqslant i \leqslant n\)）即可看出这一点。于是存在线性型 \(\psi_z\)（在 \(\mathrm{F} = \mathrm{E}_1 \otimes \mathrm{E}_2\) 上），使得

\[
\psi_ {z} (x _ {1} \otimes x _ {2}) = \overline {{\phi_ {x _ {1} , x _ {2}} (z)}} \quad (x _ {1} \in \mathrm{E} _ {1}, x _ {2} \in \mathrm{E} _ {2}).\tag{2}
\]

我们令 \(\Phi(z, t) = \psi_z(t)\)（对 F 中的 \(z, t\)）。立即可以看出 \(\Phi\) 是 \(\mathbf{E}_1 \otimes \mathbf{E}_2\) 上的一个半双线性型，它由下式刻画

\[
\Phi (x _ {1} \otimes x _ {2}, y _ {1} \otimes y _ {2}) = \langle x _ {1} | y _ {1} \rangle \langle x _ {2} | y _ {2} \rangle\tag{3}
\]

(cf.~A, IX, § 1, No.~11).

命题 1. --- 半双线性型 \(\Phi\)（在 \(\mathrm{E}_1 \otimes \mathrm{E}_2\) 上）是埃尔米特型且是正的，因而它给 \(\mathrm{E}_1 \otimes \mathrm{E}_2\) 赋予了准希尔伯特空间的结构。若 \(\mathrm{E}_1\) 与 \(\mathrm{E}_2\) 是分离的，则此空间是分离的。

公式 \(\Phi(z, t) = \overline{\Phi(t, z)}\) 当 \(z = x_1 \otimes x_2\) 且 \(t = y_1 \otimes y_2\) 时由 (3) 得出。一般情形由线性性得到，因而 \(\Phi\) 是埃尔米特型。

设 \(\mathbf{E}_1\) 与 \(\mathbf{E}_2\) 是分离的；我们将证明埃尔米特型 \(\Phi\) 是正的且非退化。设 \(z = \sum_{i=1}^{n} x_i \otimes y_i\) 为 \(\mathrm{F} = \mathrm{E}_1 \otimes \mathrm{E}_2\) 的非零元素。设 \((e_1, ..., e_n)\) 为 \(\mathbf{E}_1\) 中由 \(x_1, ..., x_n\) 生成的子空间的标准正交基（V, p.~23, 推论 1）。存在元素 \(f_1, ..., f_m\)（在 \(\mathbf{E}_2\) 中，不全为零），使得 \(z = \sum_{i=1}^{n} e_i \otimes f_i\)，因而

\[
\begin{array}{l} \Phi (z, z) = \sum_ {i, j = 1} ^ {m} \Phi (e _ {i} \otimes f _ {i}, e _ {j} \otimes f _ {j}) \\ = \sum_ {i, j} \langle e _ {i} | e _ {j} \rangle \langle f _ {i} | f _ {j} \rangle = \sum_ {i = 1} ^ {m} \| f _ {i} \| ^ {2} > 0. \end{array}
\]

对一般情形，我们现在来证明 \(\Phi\) 是正的。设 \(\tilde{E}_{i}\) 为与 \(E_{i}\) 相伴的分离准希尔伯特空间，\(\pi_{i}\) 为从 \(E_{i}\) 到 \(\tilde{E}_{i} (i = 1, 2)\) 上的典范映射。令 \(\pi = \pi_{1} \otimes \pi_{2}\)。设 \(\tilde{\Phi}\) 为 \(\tilde{E}_{1} \otimes \tilde{E}_{2}\) 上按与 \(\Phi\) 相同的方式构造的埃尔米特型。我们有

\[
\Phi (z, t) = \Phi (\pi (z), \pi (t)) \quad (z \in \mathrm{F}, t \in \mathrm{F}),
\]

而由于 \(\tilde{\Phi}\) 是正的，\(\Phi\) 也是正的。

证毕。

命题 1 中定义的准希尔伯特空间称为准希尔伯特空间 \(E_{1}\) 与 \(E_{2}\) 的张量积，记作 \(E_{1} \otimes E_{2}\)。此后我们将用 \(\langle z|t\rangle\) 表示 \(\Phi(z,t)\)，于是依定义

\[
\langle x _ {1} \otimes x _ {2} | y _ {1} \otimes y _ {2} \rangle = \langle x _ {1} | y _ {1} \rangle \langle x _ {2} | y _ {2} \rangle ;\tag{4}
\]

我们也用 \(\| z\| _2\) 或 \(\| z\|\) 来记 \(\langle z|z\rangle^{1 / 2}\)。由 (4) 得

\[
\left\| x _ {1} \otimes x _ {2} \right\| _ {2} = \left\| x _ {1} \right\|. \left\| x _ {2} \right\|,\tag{5}
\]

于是双线性映射 \((x_{1}, x_{2}) \mapsto x_{1} \otimes x_{2}\)（从 \(\mathrm{E}_1 \times \mathrm{E}_2\) 到 \(\mathrm{E}_1 \otimes \mathrm{E}_2\) 中）是连续的。

对 i = 1, 2，设 \(F_{i}\) 为 \(E_{i}\) 的向量子空间，赋以诱导的准希尔伯特结构。于是 \(F_{1} \otimes F_{2}\) 可与 \(E_{1} \otimes E_{2}\) 的一个向量子空间等同（A, II, § 7, No.~7）。公式 (4) 表明，\(F_{1} \otimes F_{2}\) 赋以由 \(E_{1} \otimes_{2} E_{2}\) 的结构诱导的准希尔伯特空间结构后，恰好就是 \(F_{1} \otimes_{2} F_{2}\)。此后我们将把 \(F_{1} \otimes_{2} F_{2}\) 与 \(E_{1} \otimes_{2} E_{2}\) 的准希尔伯特子空间等同起来。

命题 2. --- 对 \(i = 1, 2\)，设 \(\mathrm{E}_{i}\) 与 \(\mathrm{F}_{i}\) 为两个分离的准希尔伯特空间，并设 \(u_{i} \in \mathcal{L}(\mathrm{E}_{i}; \mathrm{F}_{i})\)。线性映射 \(u_{1} \otimes u_{2}\)（从 \(\mathrm{E}_{1} \otimes_{2} \mathrm{E}_{2}\) 到 \(\mathrm{F}_{1} \otimes_{2} \mathrm{F}_{2}\) 中）是连续的，并且我们有

\[
\left\| u _ {1} \otimes u _ {2} \right\| = \left\| u _ {1} \right\|. \left\| u _ {2} \right\|.
\]

考虑 \(\mathbf{E}_1\) 上由下式给出的正埃尔米特型

\[
f \left(x _ {1}, y _ {1}\right) = \| u _ {1} \| ^ {2} \left\langle x _ {1} \mid y _ {1} \right\rangle - \left\langle u _ {1} \left(x _ {1}\right) \mid u _ {1} \left(y _ {1}\right) \right\rangle .
\]

由命题 1（V, p.~25），存在正埃尔米特型 \(\Phi\)（在 \(E_1 \otimes E_2\) 上），使得

\[
\begin{array}{r l} \Phi (x _ {1} \otimes x _ {2}, y _ {1} \otimes y _ {2}) & = f (x _ {1}, y _ {1}) \langle x _ {2} | y _ {2} \rangle = \\ & = \| u _ {1} \| ^ {2} \langle x _ {1} \otimes x _ {2} | y _ {1} \otimes y _ {2} \rangle - \langle (u _ {1} \otimes 1) (x _ {1} \otimes x _ {2}) | (u _ {1} \otimes 1) (y _ {1} \otimes y _ {2}) \rangle \end{array}
\]

其中 \(x_{1}, y_{1}\) 属于 \(E_{1}\)，\(x_{2}, y_{2}\) 属于 \(E_{2}\)。由线性性得

\[
\Phi (z, t) = \| u _ {1} \| ^ {2} \langle z | t \rangle - \langle (u _ {1} \otimes 1) (z) | (u _ {1} \otimes 1) (t) \rangle
\]

其中 \(z, t\) 属于 \(\mathrm{E}_1 \otimes \mathrm{E}_2\)。由于 \(\Phi\) 是正的，得 \(\Phi(z, z) \geqslant 0\)，即 \(\| (u_1 \otimes 1).z \|_2\) \(\leqslant \| u_1\| \cdot \| z\| _2\)（对 \(z\in \mathrm{E}_1\otimes_2\mathrm{E}_2\)），亦即 \(\| u_1\otimes 1\| \leqslant \| u_1\|\)。同理可证不等式 \(\| 1\otimes u_2\| \leqslant \| u_2\|\)；又由于 \(u_{1}\otimes u_{2} = (u_{1}\otimes 1)\circ (1\otimes u_{2})\)，我们得到

\[
\left\| u _ {1} \otimes u _ {2} \right\| \leqslant \left\| u _ {1} \right\|. \left\| u _ {2} \right\|.
\]

另一方面，

\[
\begin{array}{l} \| u _ {1} \| \cdot \| u _ {2} \| = \sup _ {\| x _ {1} \| \leqslant 1, \| x _ {2} \| \leqslant 1} \| u _ {1} (x _ {1}) \| \cdot \| u _ {2} (x _ {2}) \| \\ = \sup _ {\| x _ {1} \| \leqslant 1, \| x _ {2} \| \leqslant 1} \| (u _ {1} \otimes u _ {2}) (x _ {1} \otimes x _ {2}) \| _ {2} \leqslant \| u _ {1} \otimes u _ {2} \|. \end{array}
\]

这样就完成了命题 2 的证明。

证毕。

设 \(E_1, \ldots, E_n\) 为准希尔伯特空间（\(n \geqslant 2\)）。我们用归纳法以下式定义张量积 \(E_1 \otimes_2 \ldots \otimes_2 E_n\)（也记作 \(\bigotimes_{i=1}^{n} E_i\)）

\[
\mathrm{E} _ {1} \otimes_ {2} \dots \otimes_ {2} \mathrm{E} _ {n} = (\mathrm{E} _ {1} \otimes_ {2} \dots \otimes_ {2} \mathrm{E} _ {n - 1}) \otimes_ {2} \mathrm{E} _ {n}.
\]

于是，由内积的定义，我们有

\[
\left\langle x _ {1} \otimes \dots \otimes x _ {n} \mid y _ {1} \otimes \dots \otimes y _ {n} \right\rangle = \prod_ {i = 1} ^ {n} \left\langle x _ {i} \mid y _ {i} \right\rangle ,\tag{6}
\]

特别地 \(^{1}\)

\[
\left\| x _ {1} \otimes \dots \otimes x _ {n} \right\| _ {2} = \left\| x _ {1} \right\| \dots \left\| x _ {n} \right\|,\tag{7}
\]

其中 \(x_{i},y_{i}\) 属于 \(\mathrm{E}_i\)（\(1\leqslant i\leqslant n\)）。若诸 \(\mathrm{E}_i\) 都是分离的，则 \(\mathrm{E}_1\otimes_2\dots \otimes_2\mathrm{E}_n\) 也是分离的。

设 \(F_{1}, \ldots, F_{n}\) 为准希尔伯特空间，且 \(u_{i} \in \mathcal{L}(E_{i}; F_{i})\)（\(1 \leqslant i \leqslant n\)）。对 n 用归纳法，由命题 2 推出 \(u_{1} \otimes \ldots \otimes u_{n}\) 是从 \(E_{1} \otimes_{2} \ldots \otimes_{2} E_{n}\) 到 \(F_{1} \otimes_{2} \ldots \otimes_{2} F_{n}\) 中的连续线性映射，并且

\[
\left\| u _ {1} \otimes \dots \otimes u _ {n} \right\| = \left\| u _ {1} \right\| \dots \left\| u _ {n} \right\|.\tag{8}
\]

设 \(\sigma \in \mathfrak{S}_n\) 为集合 \(\{1, 2, \dots, n\}\) 的一个置换。由 (6)，线性映射 \(p_{\sigma}\)（从 \(\mathrm{E}_1 \otimes_2 \dots \otimes_2 \mathrm{E}_n\) 到 \(\mathrm{E}_{\sigma^{-1}(1)} \otimes_2 \dots \otimes_2 \mathrm{E}_{\sigma^{-1}(n)}\) 上）由下式刻画

\[
p _ {\sigma} \left(x _ {1} \otimes \dots \otimes x _ {n}\right) = x _ {\sigma^ {- 1} (1)} \otimes \dots \otimes x _ {\sigma^ {- 1} (n)}\tag{9}
\]

是准希尔伯特空间同构（« 张量积的交换性 »）。

类似地，考虑把 \(\{1, 2, \dots , n\}\) 分成 \(m\) 个相邻区间 \(\mathrm{I}_{1}, \dots , \mathrm{I}_{m}\) 的一个划分，其中 \(\mathrm{I}_{k} = [a_{k}, a_{k + 1} - 1]\)（\(1 \leqslant k \leqslant m\)）。令

\[
\mathrm{F} _ {k} = \bigotimes_ {i = a _ {k}} ^ {a _ {k + 1} - 1} \mathrm{E} _ {i} (1 \leqslant k \leqslant m).
\]

\(^{1}\) 这里我们同样令 \(\|z\|_{2} = \langle z|z\rangle^{1/2}\)（z 在 \(E_{1} \otimes_{2} \ldots \otimes_{2} E_{n}\) 中）。

从 \(F_{1} \otimes \ldots \otimes F_{m}\) 到 \(E_{1} \otimes \ldots \otimes E_{n}\) 上的、把 \(\bigotimes_{k=1}^{m} \bigotimes_{i=a_{k}}^{a_{k+1}-1} x_{i}\) 变为 \(x_{1} \otimes \ldots \otimes x_{n}\) 的典范同构（A, II, §3, No.~9）是准希尔伯特空间同构（« 张量积的结合性 »）。

